\scanpage{119}
que c'est une \tikz[baseline=(i119.base)]{\node[inner sep=0pt] (i119) {immersion locale};
\path (i119.south west)--(i119.south east) coordinate[pos=.36] (u119a) coordinate[pos=.94] (u119b);
\draw (u119a)--(u119b);} d'espaces \deleted{\unsure{ordinaire}} \inserted{analytiques} lisses (ou lisses sur un esp. analytique donné, dans une situation relative\ldots);

\deleted{\inserted{de \unsure{codimens.|codim} $d-d'$}}
\deleted{On doit \unsure{faire|fait} l'hyp. de \unsure{finité|lissité}, que}
\deleted{Notons que pour que $D_{d,d'}\unsure{\text{---}|\longrightarrow}D_{d'}$ soit une imm. local, il suffit que}
\[
\deleted{D_{d,0}\unsure{\longrightarrow|\text{---}}D_0\quad\text{et}\quad D_{d',0}\unsure{\longrightarrow|\text{---}}D_0}
\]
% Both repeated maps below belong to the cancelled block; oblique strokes mark the cancellation.
\diagram{\begin{tikzcd}[ampersand replacement=\&,column sep=large,row sep=small,
execute at end picture={
\draw ([xshift=-3mm,yshift=5mm]r119a.north west)--([xshift=5mm,yshift=-3mm]r119e.south east);
\draw ([xshift=-8mm,yshift=8mm]r119c.north west)--([xshift=2mm,yshift=-3mm]r119d.south west);
}]
|[alias=r119a]|D_{d,0}\arrow[r,phantom,"\unsure{\longrightarrow|\text{---}}" description]\arrow[d,phantom,"\simeq"{rotate=90,description}]
\&|[alias=r119b]|D_0,\&\text{et}\&|[alias=r119c]|D_{d',0}\arrow[r,phantom,"\unsure{\longrightarrow|\text{---}}" description]\arrow[d,phantom,"\simeq"{rotate=90,description}]\&|[alias=r119d]|D_0\\
|[alias=r119e]|D_d\&\&\&D_{d'}\&
\end{tikzcd}}
\deleted{le soient, comme on voit en considérant}
% The following whole diagram is cancelled, including its arrows and both region annotations.
\diagram{\begin{tikzcd}[ampersand replacement=\&,column sep=large,row sep=large,
execute at end picture={
\path (c119b.south east)--(c119d.north west) coordinate[pos=.5] (c119comm);
\node[anchor=north,inner sep=1pt] at ([yshift=-2mm]c119comm) {$\unsure{\mathrm{comm.}|\mathrm{can.}|\mathrm{canon.}}$};
\path (c119e.north east)--(c119c.south west) coordinate[pos=.5] (c119cart);
\node[inner sep=1pt] at (c119cart) {$\mathrm{cart}$};
\draw ([xshift=-3mm,yshift=1mm]c119e.north west)--([xshift=4mm,yshift=-3mm]c119b.south east);
\draw ([xshift=-4mm,yshift=1mm]c119f.north west)--([xshift=-4mm,yshift=5mm]c119d.north west);
}]
\&|[alias=c119b]|D_{d,0}\arrow[dr]\&\\
|[alias=c119a]|D_{\unsure{d,d',0}}\arrow[ur]\arrow[r]\arrow[d,"\sim"{rotate=90,swap}]
\&|[alias=c119c]|D_{d',0}\arrow[r]\arrow[d,"\sim"{rotate=90}]\&|[alias=c119d]|D_0\\
|[alias=c119e]|D_{d,d'}\arrow[r,Rightarrow]\&|[alias=c119f]|D_{d'}\&
\end{tikzcd}}

Dans le cas lisse ou \unsure{relativ\textsuperscript{t}} lisse, l'hyp. d'imm. local équivaut d'ailleurs \unsure{:|à} celle de mono local, qui garde un sens dans tout site.

Notons que pour tt $d$ on a

\diagram{(178)\quad \begin{tikzcd}[ampersand replacement=\&,column sep=large,row sep=large]
D_{d,0}\arrow[d,"\sim"{rotate=90,swap}]\arrow[dr,-o]\&\\
D_d\&D_0\overset{\mathrm{def}}{=}X
\end{tikzcd}}
\margin{$\mathcal D$ est une catég. dans la catégorie $(\mathrm{An})_{D_0}$}
donc les $D_d$ s'envoient tous can\textsuperscript{t}. dans $X=D_0$, et on aura commutativité dans

\diagram{(179)\quad \begin{tikzcd}[ampersand replacement=\&,column sep=large,row sep=large]
D_{d,d'}\arrow[r,no head,"b",to path={(\tikztostart.east)--([xshift=-2mm]\tikztotarget.west)--++(.3mm,-.7mm)\tikztonodes}]\arrow[d,"\sigma"']\&D_{d'}\arrow[d]\\
D_d\arrow[r]\&\deleted{\unsure{D_0}}\inserted{X}
\end{tikzcd}}
\margin{118}
